\chapter*{Abstract}
\addcontentsline{toc}{chapter}{Abstract}

Sự bùng nổ của mạng lưới Internet vạn vật (IoT) đã mở rộng đáng kể bề mặt tấn công mạng, đòi hỏi các giải pháp Phát hiện Xâm nhập Mạng (NIDS) hiệu quả và bảo toàn quyền riêng tư. Học liên đoàn (Federated Learning -- FL) nổi lên như một hướng tiếp cận đầy triển vọng cho phép các thiết bị biên phối hợp huấn luyện mô hình phân tán mà không cần tập trung hóa dữ liệu lưu lượng thô. Tuy nhiên, tính phi đồng nhất về mặt thống kê (Non-IID) và hiện tượng mất cân bằng lớp cấu trúc cực đoan giữa các trạm mạng IoT biên đặt ra rào cản kỹ thuật nghiêm trọng đối với quá trình hội tụ và độ chính xác của mô hình.

Báo cáo này trình bày kết quả thực nghiệm của đề tài Khóa luận Tốt nghiệp, định vị bản chất là một \textbf{Nghiên cứu Thực nghiệm và Đo kiểm Hệ thống (Empirical Benchmarking Study)}. Nghiên cứu khảo sát hiệu năng phân loại của mô hình học máy nhẹ \texttt{TabularIoTMLP} ($14,024$ tham số, dung lượng $\approx 54.28\text{ KB}$) trên tập dữ liệu chuẩn CICIoT2023 với hơn 5.1 triệu dòng dữ liệu luồng mạng và tỷ lệ mất cân bằng lên tới $147.11:1$. Thông qua ma trận 23 kịch bản thực nghiệm đối chứng có kiểm soát (từ E1 đến E6) được phân tách thành mô hình tập trung (Centralized Baseline), dữ liệu đồng nhất (IID), chuỗi lệch phân phối Dirichlet Non-IID ($\alpha \in \{1.0, 0.5, 0.1\}$), quét toàn diện siêu tham số proximal $\mu \in \{0.001, 0.01, 0.025, 0.05\}$, và kiểm tra tính mở rộng quy mô ($K \in \{5, 7, 10\}$ clients), nghiên cứu mang lại 3 đóng góp khoa học thực chất:
\begin{enumerate}
    \item \textbf{Bằng chứng thực nghiệm phản biện (Empirical Counter-Evidence / Negative Finding):} Trong khi mô hình tập trung (E1) nhận diện tốt các lớp tấn công thiểu số (đạt tỷ lệ Recall lớp \texttt{Web-based} là $52.41\%$ và \texttt{Brute-force} là $53.36\%$), việc phân tán qua FedAvg làm tỷ lệ phát hiện \texttt{Web-based} sụp đổ xuống dưới $1.5\%$. Đặc biệt, số hạng phạt proximal trong FedProx khi $\mu \ge 0.01$ dập tắt hoàn toàn tỷ lệ phát hiện lớp thiểu số về $0.00\%$ trên mọi điều kiện phân phối (IID lẫn Non-IID). Điều này chứng minh sự thất bại của các lớp thiểu số là hệ quả toán học mang tính quy luật của phép lấy trung bình gia quyền $n_k/N$ và hiện tượng triệt tiêu tính dẻo cục bộ (Loss of Local Plasticity) của số hạng phạt proximal.
    \item \textbf{Xác lập Đường cong biên Pareto đa điều kiện (Multi-Regime Pareto Frontier):} Khảo sát thực nghiệm chứng minh $\mu = 0.001$ là điểm tối ưu Pareto phổ quát, giúp kiềm chế trôi dạt tham số từ $78.7\% - 90.8\%$ mà vẫn giải phóng tính dẻo nhận diện cho lớp tấn công hiếm \texttt{Brute-force} ($12.38\%$), đưa Macro-F1 tại mốc $\alpha=0.1$ đạt đỉnh $57.42\%$ (vượt trội hơn FedAvg $+4.17\%$ và FedProx chuẩn $+14.41\%$).
    \item \textbf{Khung đánh đổi đa chiều và Định vị Phòng thủ Chiều sâu (Defense-in-Depth):} Đo lường toàn diện chi phí truyền thông siêu nhẹ ($5.30\text{ MB} \to 10.60\text{ MB}$) tăng tuyến tính theo quy mô client với độ ổn định cao trong môi trường IID. Đồng thời, nghiên cứu xác lập ranh giới phòng thủ: tầng mạng (Layer 3/4) tối ưu cho tấn công thể tích lớn (Mirai F1 $99.70\%$, DDoS F1 $91.30\%$), trong khi tấn công Web-based cần được định tuyến cho WAF tầng ứng dụng (Layer 7).
\end{enumerate}

\textbf{Keywords:} Federated Learning, IoT Intrusion Detection, Non-IID Data, Dirichlet Distribution, FedAvg, FedProx, Pareto Frontier, CICIoT2023, Flower gRPC.

\newpage
